% ============================================================================
% §3 THE AUSTIN-LAW PROOF-SYNTHESIS BENCHMARK — v2 draft (2026-07-19)
% Drop-in replacement for the current §3. Changes from the compiled version:
%   (1) new Overview paragraph anchored on the two methodology figures;
%   (2) Problem Formulation compressed (~40% shorter, same claims);
%   (3) Corpus Construction now narrated against the figures, which carry the
%       counts; duplicated sentences removed; renewability gets the density
%       sentence (its own figure can follow);
%   (4) Answer Verification REWRITTEN to the multi-channel judge
%       (Lean = trivial side, ATP certificate = construction side), replacing
%       the stale Lean-only text; adds the descent-theorem paragraph;
%   (5) ordering: Corpus Construction now precedes Answer Verification, so the
%       section ends on "everything is machine-checked" — the note §4 opens on;
%   (6) figures get real captions and distinct labels (they currently share
%       \label{fig:methodology}, which breaks \ref).
% Citation keys marked %% VERIFY where not already in references.bib.
% ============================================================================

\section{The Austin-Law Proof-Synthesis Benchmark}

ALPS is built in three parts, displayed in Figures~\ref{fig:extension}
and~\ref{fig:funnel}: a \textbf{two-sided task} posed over laws certified to
have no nontrivial finite model, a \textbf{corpus construction} that mints such
laws at any order by extending known Austin laws and screening the results, and
an \textbf{automated judge} that verifies an answer to either side of the task
with a machine-checked certificate. We describe each component below.

\subsection{Problem Formulation}

A magma is a set $M$, the carrier, together with a single binary operation
${\op}\colon M \times M \to M$ and no further axioms. An equational law is a
universally quantified identity between two terms built from variables and
$\op$; its order is the number of times the operation appears, so the
commutative law $x \op y = y \op x$ has order two. A magma \emph{satisfies} a
law when the identity holds under every assignment of carrier elements to its
variables, and is then called a model of the law. Every ALPS instance is a law
of the form $L\colon x = T[x, y, z, \dots]$, a single variable equated with a
term $T$ of order five or more. The one-element magma is a model of every law
and so carries no information; we call it trivial, and call a model nontrivial
when its carrier holds at least two distinct elements. The question ALPS poses
for each $L$ is whether a nontrivial model exists at all: either $L$ entails
$x = y$, collapsing every model to a point, or some nontrivial magma satisfies
it.

A nontrivial model, when one exists, may be finite or infinite, and the finite
case is a matter of search: enumerating operations on carriers of increasing
size must eventually find it, requiring none of the law-specific construction
the benchmark aims to measure. We therefore admit a law only when it carries a
machine-checked proof that no nontrivial finite model exists, and call such a
law \emph{admissible}. Admissibility forces a dichotomy. Every law either
entails $x = y$ or it does not, so each instance has a determinate answer; for
an admissible law the two cases are exactly
\begin{equation}
\label{eq:dichotomy}
\underbrace{L \models x = y}_{\text{trivial}}
\quad\text{or}\quad
\underbrace{L \text{ has a nontrivial model, necessarily infinite.}}_{\text{Austin}}
\end{equation}
Laws of the second kind are the Austin laws \citep{kisielewicz1997austin}. No
instance is unanswerable: the difficulty of ALPS lies in finding the answer,
never in whether one exists.

The two sides of \eqref{eq:dichotomy} are not symmetric, and the asymmetry is
what the benchmark measures. Whether $L$ entails $x = y$ is a first-order
consequence, so a complete proof search confirms it whenever it holds; the
trivial side is, in principle, settled by deduction. The Austin side has no
such guarantee. A search that fails to derive $x = y$, however long it runs,
is not a model, and since equational entailment is undecidable in general
\citep{baader1998term}, no uniform procedure can be expected to close the gap.
Establishing that a law is Austin means producing an infinite magma that
satisfies it — a structure that is not selected from a catalogue but invented,
tailored to the individual law. This synthesis, which deductive search does
not supply, is the capability ALPS isolates.

\subsection{Corpus Construction}

\textbf{Generation by extension.} Admissible laws are too rare for random
sampling to surface in quantity, and the few it does surface are often
redundant or already enumerated. We instead grow the corpus outward from laws
already known to be Austin (Figure~\ref{fig:extension}): a seed law $x = T$ is
extended by selecting a variable $v$ in $T$ and replacing one occurrence with
$v \op w$ or $w \op v$ for some other variable $w$, yielding candidates one
order higher. Because a candidate inherits most of the seed's term structure,
it often inherits the mechanism that denies the seed a finite model, which is
why extension produces admissible laws at a rate random sampling cannot
approach. The construction is recursive — laws confirmed Austin at order $n$
seed the round at order $n{+}1$ — so the corpus has no terminal order.

\textbf{Screening.} Each candidate then passes screens ordered from cheap to
expensive (Figure~\ref{fig:funnel}). A fast proof search first removes
candidates it can already prove trivial at generation time. The remainder go
to the admissibility prover, a first-order query whose success is a proof that
no nontrivial finite model exists: when some subterm of $T$ contains every
occurrence of $x$, that subterm acts as a left inverse, making the map
$x \mapsto T$ injective; on a finite carrier an injective map is surjective,
and asserting that surjectivity lets the prover derive $x = y$. Admission
therefore requires an actual proof, and no such proof exists for a law with a
finite model, so none can enter the corpus. The price is incompleteness: when
the occurrences of $x$ are split across the root, the query certifies nothing,
and those candidates are discarded even though some are admissible. What
remains after screening is the admissible pool, split by
\eqref{eq:dichotomy} into laws proved trivial, laws whose models the screening
provers construct outright, and the residual the baseline of the next section
attacks.

\textbf{Distinctness.} Because every law descends from a small seed set,
extensions may be equivalent to their seed or to a sibling, and equivalence in
general demands a prover call per pair. For laws proved Austin, the models
found along the way decide almost all of it for free: a magma in which one law
holds and the other fails is a certificate that the two are distinct. Of the
34{,}191 pairs across the 262 laws proved Austin, 33{,}936 are separated by
such a model, 255 are proved equivalent by the prover, and none are left
undecided, collapsing the 262 laws to 195 classes — a reduction of roughly a
quarter. We report corpus size in classes rather than laws so that the count
reflects distinct problems.

\textbf{Renewability.} Extension does not exhaust itself. Normalizing by the
number of laws screened, the yield of Austin laws \emph{rises} with order —
20.5, 24.4, and 25.4 per thousand screened at orders six, seven, and eight —
and each order contributes classes no earlier order is equivalent to.
%% (density figure can anchor this sentence when built)
Since the construction applies at every order, the corpus can always be
extended past whatever has been enumerated or released before. Evaluation can
therefore always draw on instances minted after the model under test was
trained: the corpus is renewable, and in turn contamination-free, in a way no
fixed problem set can be.

\subsection{Answer Verification}

A solved ALPS instance is a verified mathematical fact, not a predicted label.
The two sides of the task call for different machinery, so the judge issues a
certificate through the channel that fits the answer: a proof checked by the
Lean kernel for the trivial side, and an automated-prover certificate for a
constructed model. Both channels are fully automatic, and a submission is
accepted only when its check succeeds.

\textbf{The trivial channel.} For each law the judge generates the proposition
\begin{equation}
\label{eq:trivialgoal}
\textsc{TrivialGoal} \equiv \forall M\ \forall {\op}\colon M \times M \to M,\;
\mathrm{Law}(\op) \rightarrow \forall a, b : M,\ a = b,
\end{equation}
where $\mathrm{Law}(\op)$ asserts that $\op$ satisfies $L$, and a submission
is a Lean proof of it. The statement is generated, never written by the
submitter: the judged file sandwiches the submitted proof between a header
that fixes \eqref{eq:trivialgoal} and a footer that forces the proof to
elaborate at exactly that type, and the proof's axiom footprint must fall
within a fixed allowlist. A proof the kernel accepts under these constraints
is correct by construction \citep{moura2021lean}.

\textbf{The construction channel.} A solver that believes $L$ is Austin does
not write a Lean proof about an infinite object. It submits the model itself,
as a finite presentation $E$ — a set of equations over $\op$ that pins down
the intended structure. The judge then runs two independent prover queries
\citep{kovacs2013first}:
\begin{equation}
\label{eq:construction}
\text{(a)}\;\; E \vdash L
\qquad\qquad
\text{(b)}\;\; E \cup \{\,a \neq b\,\} \text{ is satisfiable,}
\end{equation}
certifying the submission only when both succeed. Query (a) is a refutation
proof that every model of $E$ satisfies $L$; query (b) is a terminating
saturation establishing that $E$ has a model with two distinct elements.
Together they exhibit a nontrivial model of $L$, which admissibility forces to
be infinite, so $L$ is Austin. The pair is also self-policing. Submitting the
law itself as $E$ gains nothing: (a) becomes immediate, but (b) becomes the
bare saturation of $L$, which is exactly what diverges on hard instances. A
collapsing presentation such as $\{x = y\}$ passes (a) but fails (b). The
judge accepts a certificate from any trusted checker — the construction
certificates are cross-checkable by a second prover
\citep{smallbone2021twee}, and a Lean proof of the corresponding existential
goal would equally count — so its trust base is stated plainly: the Lean
kernel on the trivial side, prover saturation on the construction side, the
same base on which the corpus labels themselves rest.

\textbf{Why two channels.} A single Lean arbiter for both sides is not
available, for a reason that strengthens the benchmark's premise. If an
admissible law had a model whose operation is a polynomial or affine formula
over $\mathbb{Z}$, $\mathbb{Z}^k$, or $\mathbb{Z}/n$, the identity would
survive reduction modulo two and yield a nontrivial \emph{finite} model,
contradicting admissibility. No Austin model is an arithmetic formula: the
familiar, easily formalized model families are provably empty here, so an
Austin answer cannot be a memorized pattern, and the structures that do work —
term models built from saturations — await a formalization of ground
confluence that we develop as companion work. The two-channel judge is what
makes the construction side gradable today, at the cost of one additional
trusted checker.

% ---------------------------------------------------------------------------
% FIGURE CAPTIONS (replace the placeholders)
% ---------------------------------------------------------------------------
% Extension engine (currently Figures/ALPS_Methodology_a.pdf):
%
% \caption{Corpus generation by extension. A seed Austin law is extended by
% replacing one variable occurrence $v$ with $v \op w$, yielding candidate
% laws one order higher that tend to inherit the seed's lack of finite models.
% Laws confirmed Austin re-seed the next round, so generation applies
% recursively at every order.}
% \label{fig:extension}
%
% Screening funnel (currently Figures/ALPS_Methodology_b.pdf):
%
% \caption{Screening the candidate pool. Each candidate passes screens ordered
% from cheap to expensive: a finite-model filter discards laws with small
% nontrivial models, and the admissibility prover certifies the rest or
% discards what it cannot certify. Band width tracks the number of surviving
% laws; the admissible pool splits into laws proved trivial, Austin classes
% with exhibited models, and the unresolved hard tier. Counts are provisional
% pending the full portfolio sweep.}
% \label{fig:funnel}
% ---------------------------------------------------------------------------
